\documentclass[12pt]{article}
\usepackage[utf8]{inputenc}
\usepackage[T1]{fontenc}
\usepackage{hyperref}
\title{Puruborá Language Resources for K-12 Curriculum Development}
\author{Compiled by Claude Sonnet 5.5 (claude-sonnet-5-5) for Andr\'as Kornai}
\date{}
\begin{document}
\maketitle

\section{Introduction}
Puruborá (ISO 639-3 \texttt{pur}, Glottocode puru1264) is a nearly extinct Tupian language of Rondônia, Brazil (Costa Marques and the upper Rio São Miguel). Search reports about two remaining speakers who recall mainly words and phrases. Documentation exists as phonological studies and a lexical database of about 700 items (Museu Goeldi), but no grammar or school materials were found.

\section{Primary Educational Resources}
None found. A planned Puruborá--Portuguese vocabulary is mentioned in the Goeldi record, but no published edition was seen.

\section{Linguistic and Reference Materials}
\begin{itemize}
\item Purubora language (Wikipedia): \url{https://en.wikipedia.org/wiki/Purubora_language}
\item Overview of Puruborá language characteristics (Portuguese, Povos Ind\'igenas no Brasil): \url{https://pib.socioambiental.org/files/file/PIB_verbetes/purubor%c3%a1/leiamais_caracteristicas%20da%20lingua%20purubora.pdf}
\item Museu Goeldi repository record, Puruborá documentation: \url{https://repositorio.museu-goeldi.br/handle/mgoeldi/2304}
\item Phonetic and phonological analysis of Puruborá (PDF, Museu Goeldi): \url{https://repositorio.museu-goeldi.br/bitstream/mgoeldi/2311/1/70-AN%c3%81LISE%20FON%c3%89TICA%20E%20FONOL%c3%93GICA%20DA%20L%c3%adNGUA%20PURUBOR%c3%81.pdf}
\end{itemize}

\section{Religious and Community Resources}
None found.

\section{Digital and Audio Resources}
\begin{itemize}
\item Endangered Languages Project entry: \url{https://www.endangeredlanguages.com/lang/2766}
\end{itemize}
Wikipedia links an ELAR archive of Tupian languages including Puruborá, but its URL was not seen in search results and is not listed.

\section{Curriculum Development Context}
With about two speakers and no pedagogical materials, any programme would be revitalization work built on the existing lexical database and phonological analyses, led by the Puruborá community.

\end{document}